%========================================================
% CAPÍTULO 6 — RESULTADOS E DISCUSSÃO
% Incluído via \input em main.tex. NÃO adicionar preâmbulo.
% Separar sempre resultado CONFIRMADO (medido e sincronizado) de [a coletar].
%========================================================
\section{Resultados e discussão}\label{cap:resultados}

Este capítulo confronta os critérios de sucesso fixados na Seção~\ref{sec:protocolo}
com o que foi efetivamente medido no protótipo DOMUS. Adota-se, por rigor, uma
distinção explícita entre \textit{resultados confirmados} --- obtidos por medição
instrumentada e verificados em hardware ou no banco de dados do \textit{hub} --- e
métricas \textit{em coleta}, cujo procedimento está definido, porém cuja aquisição
ainda estava em andamento no fechamento desta versão. O Quadro~\ref{tab:panorama}
resume esse estado. A opção por declarar o que ainda não foi medido, em vez de
preencher lacunas com estimativas, é coerente com a postura de honestidade documental
assumida no trabalho e evita que a discussão confunda promessa de projeto com
evidência empírica. Vale registrar que as métricas de voz confirmadas não permanecem
apenas no banco: o próprio sistema as expõe em um painel de \textit{Resultados} que as
recalcula ao vivo a partir da tabela \texttt{voice\_commands} e distingue visualmente o
que é medido do que está em coleta --- transportando para dentro do produto a mesma
honestidade documental adotada aqui.

\begin{table}[H]
\centering
\caption{Panorama dos critérios de avaliação e seu estado de coleta.}
\label{tab:panorama}
\begin{tabular}{@{}lll@{}}
\toprule
\textbf{Critério} & \textbf{Instrumento} & \textbf{Estado} \\
\midrule
Controle local de hardware   & Latência KLAP (Tapo P110)        & Confirmado \\
Acurácia de intenção de voz  & \textit{Log} \texttt{voice\_commands} & Confirmado \\
Sucesso de execução de voz   & \textit{Log} \texttt{voice\_commands} & Confirmado \\
Latência voz$\rightarrow$ação & Percentis p50/p95                & Confirmado \\
Taxa de erro de palavra (WER) & Corpus controlado                & A coletar  \\
Custo total (BOM)            & Planilha datada, 2 cenários      & A coletar  \\
Usabilidade (SUS)            & Questionário, $n=3$--$5$         & A coletar  \\
Economia de energia          & Série antes/depois com carga     & A coletar  \\
Comissionamento sem \textit{app} & Tapo (pleno) / EWS 410 (bloqueado) & Confirmado \\
\bottomrule
\end{tabular}
\end{table}

\subsection{Controle local de hardware}\label{sec:res-controle}

O critério mais elementar da tese --- controlar hardware de prateleira sem trânsito
obrigatório pela nuvem do fabricante --- foi \textbf{confirmado} para a tomada
inteligente TP-Link Tapo~P110. O adaptador correspondente estabelece a sessão
diretamente na rede local pelo protocolo proprietário KLAP (\textit{Key Local
Authentication Protocol}), negociando o \textit{handshake} de autenticação a partir de
credenciais de consumidor, sem recorrer aos servidores da TP-Link em regime permanente
(TP-LINK, 2024). Trata-se da materialização operacional do princípio \textit{local-first}
descrito por Kleppmann et~al.\ (2019): o plano de controle reside na instalação local.

As latências fim-a-fim das operações de comando, medidas na rede local do laboratório,
situaram-se na faixa de \textbf{201 a 332~ms} (Tabela~\ref{tab:latencia-tapo}). Esse
intervalo é imperceptível como atraso de interação para o usuário e comparável ao de
comandos que trafegariam pela nuvem em condição de rede favorável, com a vantagem de
não depender da conectividade externa --- justamente o gargalo estrutural identificado
na justificativa do trabalho, dado que apenas 22\% dos domicílios brasileiros dispõem
de conectividade satisfatória (CETIC.br, 2024).

\begin{table}[H]
\centering
\caption{Latências de controle local do Tapo P110 pelo protocolo KLAP (rede local).}
\label{tab:latencia-tapo}
\begin{tabular}{@{}lcc@{}}
\toprule
\textbf{Estatística} & \textbf{Latência (ms)} & \textbf{Observação} \\
\midrule
Menor latência observada  & 201   & Comando após sessão estabelecida \\
Maior latência observada  & 332   & Inclui renegociação de \textit{handshake} \\
Mediana (p50)             & [...] & A confirmar no \textit{log} \\
Amostras ($n$)            & [...] & A confirmar no \textit{log} \\
\bottomrule
\end{tabular}
\end{table}
% NOTA: a faixa 201–332 ms está CONFIRMADA. Preencher p50, n e, se possível, quebra por
% operação (ligar / desligar / leitura de energia) a partir dos tempos registrados nos
% adaptadores. Se houver série suficiente, considerar figura fig_latencia_tapo.pdf.

\subsection{Reconhecimento de voz}\label{sec:res-voz}

O \textit{pipeline} de voz foi avaliado sobre \textbf{138 comandos reais} registrados na
tabela \texttt{voice\_commands} durante o uso do protótipo, o que caracteriza uma
coleta em uso, e não em laboratório controlado. A Tabela~\ref{tab:voz} consolida as
métricas confirmadas. É central, para a leitura correta dos números, separar quatro
grandezas que a literatura e o senso comum tendem a fundir: a acurácia de reconhecimento
da fala (WER), a acurácia de \textit{intenção}, o sucesso de \textit{execução} e a
latência.

A \textbf{acurácia de intenção} --- fração dos comandos em que o \textit{parser}
extraiu corretamente a ação pretendida e seu alvo --- foi de \textbf{87,7\%}. O
resultado indica que o interpretador de linguagem natural em pt-BR desempenha o seu
papel de forma robusta mesmo sobre transcrições ruidosas. Já o \textbf{sucesso de
execução} --- fração dos comandos que resultaram na mudança de estado física
pretendida --- foi de apenas \textbf{66,7\%}. A diferença de mais de vinte pontos
percentuais entre as duas métricas \emph{não} decorre de falha do \textit{parser}: ela
é \textbf{dominada por hardware indisponível} no momento do comando (dispositivo
\textit{offline}, tomada sem energia, alvo fora da rede). Em outras palavras, o sistema
compreendeu a intenção, mas não pôde efetivá-la porque o dispositivo-alvo não estava
acessível --- uma falha de disponibilidade de hardware, não de reconhecimento. Essa
decomposição é a resposta direta à objeção previsível da banca de que ``a execução
ficou abaixo de dois terços'': o número mede a confiabilidade do parque físico
disponível durante a coleta, não a qualidade do interpretador.

\begin{table}[H]
\centering
\caption{Métricas confirmadas do \textit{pipeline} de voz ($n=138$ comandos reais).}
\label{tab:voz}
\begin{tabular}{@{}lcl@{}}
\toprule
\textbf{Métrica} & \textbf{Valor} & \textbf{Observação} \\
\midrule
Comandos instrumentados ($n$) & 138      & Tabela \texttt{voice\_commands} \\
Acurácia de intenção          & 87,7\%   & Intenção + alvo corretos \\
Sucesso de execução           & 66,7\%   & Limitado por hardware \textit{offline} \\
Latência p50                  & 463~ms   & Voz $\rightarrow$ ação \\
Latência p95                  & 2.140~ms & Cauda superior \\
Latência máxima (\textit{outlier}) & 12,3~s & \textit{Timeout} de hardware \\
WER (taxa de erro de palavra) & [...]    & A coletar (corpus controlado) \\
\bottomrule
\end{tabular}
\end{table}

Quanto à \textbf{latência voz$\rightarrow$ação}, a mediana (p50) de \textbf{463~ms}
confirma uma resposta perceptualmente imediata no caso típico. O percentil~95 de
\textbf{2.140~ms} revela uma cauda superior mais lenta, apenas ligeiramente acima do
patamar informal de dois segundos e explicável pelo custo de transcrição do
\textit{Whisper} no \textit{hub} somado ao tempo de resposta do dispositivo. O
\textbf{\textit{outlier} máximo de 12,3~s} não representa lentidão do reconhecimento:
corresponde a um \textbf{\textit{timeout} de hardware}, em que o adaptador aguardou até
o limite a resposta de um dispositivo que não respondeu --- mais uma manifestação do
mesmo gargalo de disponibilidade física discutido acima. Reportar a latência por
percentis, e não por média, é deliberado: a média seria contaminada por esses eventos
de \textit{timeout} e mascararia o comportamento real percebido pela maioria dos
comandos.

Um achado qualitativo merece registro por seu valor probatório espontâneo. Em pelo
menos um enunciado, o \textit{Whisper} produziu a transcrição alucinada
\texttt{"[Som de futebol]"} para um segmento sem fala inteligível. Longe de ser um
constrangimento a ocultar, o episódio é \textbf{evidência empírica de uma limitação
conhecida} do reconhecimento de fala de vocabulário aberto em larga escala: modelos do
tipo \textit{Whisper} tendem a ``alucinar'' rótulos plausíveis em trechos de silêncio
ou ruído (RADFORD et~al., 2022). O fato de a alucinação ter emergido no uso real, e não
em um teste adversarial construído, reforça a fidelidade da coleta e alimenta
diretamente a discussão de limitações e de robustez da demonstração (Seção~\ref{sec:res-comissionamento}
e Capítulo~\ref{cap:consideracoes}).

A \textbf{taxa de erro de palavra (WER)} permanece \textbf{a coletar}. Diferentemente
das métricas acima, ela exige um corpus controlado --- protocolo previsto de, no mínimo,
50 enunciados com pelo menos três falantes distintos e transcrição de referência ---,
acompanhado de uma matriz de confusão de intenções. Sem esse corpus, seria incorreto
reportar WER a partir dos 138 comandos de uso, que não possuem \textit{ground truth}
textual.
% INSERIR (semana de métricas): WER global e por falante + matriz de confusão de
% intenções. Considerar figura fig_matriz_confusao.pdf.

\begin{figure}[H]
\centering
\includegraphics[width=0.82\textwidth]{05-resultados.png}
\caption{Painel de \textit{Resultados} do DOMUS: as métricas de voz são recalculadas ao
vivo a partir da tabela \texttt{voice\_commands}, distinguindo visualmente o que é
\emph{medido} do que está \emph{em coleta} --- a mesma honestidade documental desta seção.}
\label{fig:painel-resultados}
\end{figure}

\subsection{Custo (BOM)}\label{sec:res-custo}

O custo total de materiais (\textit{Bill of Materials}, BOM) é a métrica que sustenta
diretamente a alegação de ``baixo custo'' do título. Seu levantamento definitivo está
\textbf{a coletar}: exige uma planilha única e datada, com preços de varejo referenciados
(notas fiscais ou \textit{links}), estruturada em \textbf{dois cenários} que enfrentam
explicitamente a questão do \textit{hub} --- o item que a narrativa ingênua de custo
tende a esconder. A Tabela~\ref{tab:bom} apresenta o esqueleto a ser preenchido.

\begin{table}[H]
\centering
\caption{Custo de materiais (BOM) do DOMUS em dois cenários --- valores a coletar.}
\label{tab:bom}
\begin{tabular}{@{}lcc@{}}
\toprule
\textbf{Item} & \textbf{Incremental (R\$)} & \textbf{Completo (R\$)} \\
\midrule
TP-Link Tapo P110 (tomada)          & [...] & [...] \\
Intelbras EWS 410 (lâmpada)         & [...] & [...] \\
\textit{Hub} de processamento       & 0\,\textsuperscript{a} & [...]\,\textsuperscript{b} \\
Outros (cabos, fonte)               & [...] & [...] \\
\midrule
\textbf{Total}                      & \textbf{[...]} & \textbf{[...]} \\
\bottomrule
\end{tabular}

{\footnotesize\noindent\textsuperscript{a}~PC/notebook doméstico reaproveitado como
\textit{hub} (custo marginal zero). \quad
\textsuperscript{b}~Mini-PC ou Raspberry~Pi dedicado.\par}
\end{table}
% INSERIR os valores datados; discutir a meta informal de R$ 200 de custo INCREMENTAL,
% e o fato de que a divergência histórica (R$ 130 vs R$ 180 na documentação) se resolve
% justamente por explicitar se o hub entra ou não na conta.

A discussão central desta subseção, portanto, não é o preço da lâmpada ou da tomada
--- baratos e já premissa do trabalho ---, mas o \textbf{custo do \textit{hub}}. No
cenário \textit{incremental}, em que um computador doméstico ocioso é reaproveitado, o
custo marginal do \textit{hub} tende a zero e a meta de custo total incremental
inferior a R\$~200 é plausível. No cenário \textit{completo}, com um mini-PC ou
Raspberry~Pi dedicado, esse custo precisa ser contabilizado com honestidade, sob pena
de a alegação de democratização escamotear o seu principal item de despesa.

\subsection{Usabilidade (SUS)}\label{sec:res-sus}

A avaliação de usabilidade pela \textit{System Usability Scale} (SUS) está \textbf{a
coletar}. O protocolo previsto aplica o questionário SUS padrão a uma amostra de
$n = 3$ a $5$ participantes, após um roteiro de cinco tarefas representativas
(comissionar um dispositivo, emitir um comando de voz, criar uma automação por horário,
consultar o painel de energia e recuperar-se de uma falha simulada). A
Tabela~\ref{tab:sus} apresenta a estrutura de reporte.

\begin{table}[H]
\centering
\caption{Resultado da avaliação de usabilidade (SUS) --- a coletar.}
\label{tab:sus}
\begin{tabular}{@{}lc@{}}
\toprule
\textbf{Item} & \textbf{Valor} \\
\midrule
Participantes ($n$)          & [...] \\
Pontuação SUS média          & [...] \\
Desvio-padrão                & [...] \\
Faixa (mín.\ --- máx.)       & [...] \\
Interpretação qualitativa    & [...] \\
\bottomrule
\end{tabular}
\end{table}
% INSERIR: SUS por participante + média. Score >= 68 é a média de referência da escala;
% discutir com cautela dado o n pequeno. Reconhecer explicitamente que n=3–5 sustenta
% um argumento DIRECIONAL de usabilidade, não uma generalização — resposta preparada
% para a objeção de "democratização com n=1".

Reconhece-se, desde já, que uma amostra dessa magnitude não autoriza generalização
estatística. O propósito da medição é \textbf{direcional}: verificar se usuários sem
formação técnica conseguem completar as tarefas essenciais e identificar barreiras
grosseiras de usabilidade, alimentando a discussão de democratização em sua forma
prudente. Sem coleta de SUS, a defesa deve recuar para o título alternativo mais
conservador, centrado em viabilidade técnica.

\subsection{Energia}\label{sec:res-energia}

O painel de energia foi implementado e alimentado pela telemetria do Tapo~P110, que o
\texttt{EnergyService} coleta em intervalos regulares e persiste na tabela
\texttt{energy\_readings}, acumulando um volume expressivo de leituras. Sobre essa série,
o painel agrega o consumo \textbf{da casa inteira} --- somando todas as conexões que
medem energia --- e o apresenta em janelas de \textbf{24~h, 7 e 30~dias}, com resolução
que desce ao \textbf{minuto} na janela diária; complementa a leitura com a
\textbf{participação de cada conexão} no total e um \textbf{comparativo mês a mês}, além
do consumo por aparelho acessível ao abrir cada dispositivo.

\begin{figure}[H]
\centering
\includegraphics[width=0.82\textwidth]{04-energia.png}
\caption{Painel de energia da casa: consumo agregado com seletor de período (24~h em
resolução de minutos, 7 e 30 dias), participação por conexão e comparativo mensal,
alimentado pela telemetria real do Tapo~P110.}
\label{fig:painel-energia}
\end{figure}

Contudo, a avaliação de
\textbf{economia de energia} propriamente dita está \textbf{a coletar}, por uma razão
metodológica honesta: as leituras acumuladas foram obtidas com a tomada sem carga útil
acoplada (potência residual da ordem de fração de \textit{watt}), o que torna a série
inadequada para demonstrar economia real. A avaliação válida requer uma \textbf{série
temporal antes/depois com carga real} (por exemplo, um ventilador ou abajur submetido a
uma rotina de desligamento programado), coletada por um período mínimo de vários dias
consecutivos. A Tabela~\ref{tab:energia} apresenta a estrutura de reporte.

\begin{table}[H]
\centering
\caption{Consumo antes/depois da automação, com carga real --- a coletar.}
\label{tab:energia}
\begin{tabular}{@{}lccc@{}}
\toprule
\textbf{Condição} & \textbf{Período} & \textbf{Consumo (kWh)} & \textbf{Custo (R\$)} \\
\midrule
Antes (sem automação) & [...] & [...] & [...] \\
Depois (com automação) & [...] & [...] & [...] \\
\midrule
\textbf{Economia}     & --- & \textbf{[...]} & \textbf{[...]} \\
\bottomrule
\end{tabular}
\end{table}
% INSERIR: série >= 5–7 dias com carga real; tarifa em R$/kWh referenciada (não
% hardcoded). Se a série não for viável a tempo, REBAIXAR explicitamente para "painel
% demonstrado; economia não avaliada" nas limitações (§7.1), sem inflar o resultado.

Enquanto a série com carga não estiver disponível, o resultado honesto é: o
\emph{painel} de monitoramento de energia está \textbf{demonstrado e funcional}, porém a
\emph{economia} proporcionada pela automação \textbf{não foi avaliada}. Essa distinção
é mantida deliberadamente para não converter capacidade de medição em alegação de
benefício não comprovado.

\subsection{Validação com usuários}\label{sec:res-survey}

Para confrontar as premissas do trabalho --- sobretudo a de que a barreira dominante é o
\emph{comissionamento}, e não o preço --- com a percepção de usuários reais, aplicou-se um
questionário eletrônico (Apêndice~\ref{ap:survey}) a uma amostra de conveniência de
diferentes faixas etárias e regiões (Recife/PE, São Paulo, Rio de Janeiro e participantes
na Itália). Trata-se de uma avaliação \textbf{direcional} ($n = 8$--$15$), não de uma
inferência estatística. No fechamento desta versão, a coleta estava \textbf{em andamento}.
Os resultados consolidados serão reportados aqui em gráficos gerados a partir das respostas
reais: (i)~perfil etário e regional da amostra; (ii)~fração que já possui automação
\emph{versus} a que automatizaria se fosse simples e barata; e, sobretudo, (iii)~o
\emph{ranking das dificuldades percebidas}, cuja leitura será cruzada com o achado de
comissionamento (Seção~\ref{sec:res-comissionamento}) para testar, com evidência de
usuários, a hipótese central do trabalho.

% INSERIR (após a coleta no Google Forms): gerar dos dados REAIS e NÃO antes:
%  (a) barras horizontais do ranking de barreiras (Bloco 3 do instrumento);
%  (b) perfil etário/regional + % que já tem automação vs. automatizaria;
%  (c) SUS médio por faixa etária/cidade. Não reportar números sem coleta real.

\subsection{Comissionamento}\label{sec:res-comissionamento}

O comissionamento --- transformar hardware de prateleira em dispositivo controlável
sem recorrer ao aplicativo proprietário nem a portais de desenvolvedor --- é a
contribuição-título do trabalho, e os resultados aqui são \textbf{assimétricos e, por
isso mesmo, esclarecedores}.

No caminho \textbf{TP-Link Tapo}, a forma plena foi \textbf{demonstrada}: partindo de
credenciais de consumidor, o DOMUS estabelece controle local pelo protocolo KLAP
(Seção~\ref{sec:res-controle}), sem que o uso cotidiano dependa do \textit{app} do
fabricante ou de servidores externos (TP-LINK, 2024). Esse é o caso que sustenta a
forma \textit{forte} da tese em um dispositivo real.

No caminho \textbf{Intelbras/Tuya EWS 410}, o comissionamento foi \textbf{bloqueado no
provisionamento}. A lâmpada nunca chegou a ser controlada em hardware porque o fluxo
exigiu, na prática, credenciais de baixo nível (\textit{Device~ID} e
\textit{local\_key}) obtidas por meio do portal de desenvolvedor Tuya~IoT, cujo
emparelhamento de dispositivos esbarrou no \textbf{limite de dois \textit{slots} de
dispositivo} da plataforma em regime de avaliação (TUYA, 2025). Um \textit{harness} de
provisionamento local foi preparado, mas o gargalo permaneceu no ato de vincular o
dispositivo à conta da plataforma.

A leitura deste bloqueio é o ponto mais importante do capítulo. Ele \textbf{não é um
fracasso de engenharia a ser minimizado}: é a \textbf{evidência empírica mais forte da
própria tese}. A barreira de comissionamento que a Introdução postula como o real
obstáculo à democratização --- aplicativos proprietários, contas em nuvem, portais de
desenvolvedor em inglês e limites administrativos de plataforma --- foi \textbf{real a
ponto de bloquear até o autor do trabalho}, um desenvolvedor com acesso ao \textit{harness},
à documentação e ao ferramental. Se essa barreira detém quem a estuda, seu efeito sobre
o domicílio típico é qualitativamente mais severo, o que é consistente com a baixa
penetração de automação residencial no país (IBGE, 2024). O contraste entre o caminho
Tapo (pleno) e o caminho Tuya (bloqueado) converte-se, assim, de limitação em resultado:
demonstra simultaneamente que o objetivo é \textit{alcançável} (Tapo) e que a barreira
que ele ataca é \textit{concreta e mensurável} (Tuya). A validação plena do caminho Tuya
permanece em andamento e é retomada como limitação declarada no
Capítulo~\ref{cap:consideracoes}.
% NOTA: a promessa "sem app / sem portal" está DEMONSTRADA para Tapo e CARACTERIZADA
% como gap para Tuya. Não afirmar controle da EWS 410 em hardware. Se o timebox da EWS
% 410 for bem-sucedido antes da defesa, atualizar esta subseção para a forma forte.
